\documentclass[10.5pt]{ctexart}
\usepackage[a4paper,top=1.45cm,bottom=1.45cm,left=1.55cm,right=1.55cm]{geometry}
\usepackage{hyperref}
\usepackage{enumitem}
\usepackage{multicol}
\usepackage{titlesec}
\usepackage{xcolor}
\usepackage{fancyhdr}
\setlength{\parindent}{0pt}
\setlength{\parskip}{4pt}
\setlist{nosep,leftmargin=1.3em}
\definecolor{ormuted}{HTML}{626262}
\definecolor{orline}{HTML}{B9B5AE}
\hypersetup{colorlinks=true,urlcolor=black,linkcolor=black}
\pagestyle{fancy}
\fancyhf{}
\lhead{\small Opportunity Radar · OR Morning · 2026-10-13}
\rhead{\small\thepage/8}
\cfoot{}
\renewcommand{\headrulewidth}{0.35pt}
\titleformat{\section}{\Huge\bfseries}{ }{0pt}{}
\begin{document}
{\small\bfseries OR MORNING / SECTION 01 / 2026-10-13}
\vspace{2mm}
\hrule height 0.9pt
\vspace{1.5mm}
\hrule height 0.25pt
\vspace{4mm}
\section*{今日总览}
{\color{ormuted} 如果今天只能看三分钟，这一版必须已经足够。}
\vspace{2mm}
{\fontsize{25}{29}\selectfont\bfseries 行政机构\_中国人民大学高瓴人工智能学院}\par
\vspace{2mm}
{\large 本版筛出 7 个高信号事件。}
\vspace{3mm}
{\Large\bfseries 行政机构\_中国人民大学高瓴人工智能学院}\par
首页 学院概况 学院简介 数说学院 现任领导 学术委员会 行政机构 新闻公告 学院新闻 学院公告 师生风采 光影视界 师资队伍 专任教师 博士后队伍 人才招聘 人才培养 本科生 学术型硕士研究生 专业学位硕士研究生 学术型博士研究生 专业学位博士研究生 在职课程培训班 同等学力申请硕士学位 科学研究 研究机构 研究成果 合作交流 对外合作 国际交流 校友名录 光影记录 校友名录 高礼研究院 EN 教师入口 学院概况 学院简介 数说学院 现任领导 学术委员会 行政机构 行政机构 您当前所在的位置： 首页 > 学院概况 > 行政机构 教务工作 - 教务管理：负责排课、选课、成绩管理，以及毕业实践和毕业论文的组织与管理，包括出具成绩单和协助学历研究生招生工作。 - 教师教学工作管理：计算和考核教师的教学工作量，确保教学任务的合理分配和履行。 - 教学委员会支持：协助各教学工作委员会的秘书处工作，包括会议组织、文件管理等。 - 学科建设
\textbf{与你有关：} 这是当前用户可实际参与或申请的高可触达机会；参与范围：北京 / 人大；命中用户订阅版面：opportunity、research、ai；来源接近一手且权威度高；标题/正文包含高信号变化
\textbf{建议：} 检查报名资格、截止时间、参与形式和预期收益；合适则立即加入任务系统。
{\footnotesize\color{ormuted} 来源：中国人民大学高瓴人工智能学院公告 · \url{https://ai.ruc.edu.cn/xygk/xzjg/index.htm}}
\vspace{3mm}
\hrule height 0.25pt
\vspace{3mm}
\begin{multicols}{2}
{\large\bfseries 姚班培养-清华大学交叉信息研究院}\par
首页 热点新闻 讲座预告 院系概况 院系简介 院长寄语 现任领导 院系沿革 新闻动态 院系动态 媒体聚焦 人员队伍 全职教师 研究系列 荣誉客座 兼职教师 博士后 行政人员 实验员 姚班本科生 姚班培养 姚班招生 研究生 研究生培养 研究生招生 量子中心 中心简介 师资力量 科研概况 量子中心动态 科学研究 组织架构 课题组介绍 研究动态 科研简报 校园生活 学生活动 衣食住行 公开招聘 教师招聘 博士后招聘 职员招聘 校友频道 校友风采 校友俱乐部 校友捐赠 姚班本科生 姚班培养 姚班招生 清华学堂计算机科学实验班 姚班介绍 “清华学堂计算机科学实验班（姚班）”由世界著名计算机科学家姚期智院士于2005年创办，是清华大学培养国际顶尖计算机科学及交叉创新人才的本科生实验班，被誉为“最优秀的本科生和最优秀的本科教育”标杆。 姚班面向科技发展前沿与国家重大需求构建培养体系，以计算机科学、人工智能和量子信息三大前沿专业方向为核心，供学
\textbf{与你有关：} 这是当前用户可实际参与或申请的高可触达机会；参与范围：北京 / 清华；命中用户订阅版面：opportunity、research、ai；来源接近一手且权威度高；标题/正文包含高信号变化
\textbf{建议：} 检查报名资格、截止时间、参与形式和预期收益；合适则立即加入任务系统。
{\footnotesize\color{ormuted} 来源：清华交叉信息研究院 · \url{https://iiis.tsinghua.edu.cn/ybbks/ybpy.htm}}
\vspace{2mm}
\hrule height 0.2pt
\vspace{2mm}
{\large\bfseries 中心简介-清华大学交叉信息研究院}\par
首页 热点新闻 讲座预告 院系概况 院系简介 院长寄语 现任领导 院系沿革 新闻动态 院系动态 媒体聚焦 人员队伍 全职教师 研究系列 荣誉客座 兼职教师 博士后 行政人员 实验员 姚班本科生 姚班培养 姚班招生 研究生 研究生培养 研究生招生 量子中心 中心简介 师资力量 科研概况 量子中心动态 科学研究 组织架构 课题组介绍 研究动态 科研简报 校园生活 学生活动 衣食住行 公开招聘 教师招聘 博士后招聘 职员招聘 校友频道 校友风采 校友俱乐部 校友捐赠 量子中心 中心简介 师资力量 科研概况 量子中心动态 量子中心 中心介绍 量子信息是一门新兴的交叉学科领域，涉及计算机科学、信息科学和量子物理，为后摩尔时代提供了新的技术路径，探索了21世纪信息科学发展的新基础与新方法。量子信息技术已成为各国战略竞争的焦点之一。量子信息科学作为量子调控的核心，是信息网络与纳米科学等领域的交叉方向，代表了中国未来科技发展的重要战略方向之一。
\textbf{与你有关：} 这是当前用户可实际参与或申请的高可触达机会；参与范围：北京 / 清华；命中用户订阅版面：opportunity、research、ai；来源接近一手且权威度高；标题/正文包含高信号变化
\textbf{建议：} 检查报名资格、截止时间、参与形式和预期收益；合适则立即加入任务系统。
{\footnotesize\color{ormuted} 来源：清华交叉信息研究院 · \url{https://iiis.tsinghua.edu.cn/lzzx/zxjj.htm}}
\vspace{2mm}
\hrule height 0.2pt
\vspace{2mm}
{\large\bfseries 最新日程：清华大学第二十一届公共管理前沿论坛即将开幕}\par
最新日程：清华大学第二十一届公共管理前沿论坛即将开幕 IR AI SoftEng 清华大学人工智能国际治理研究院 2026-10-10 22:50 分享 以下文章来源于mp.weixin.qq.com 点击蓝字 关注我们 清华大学第二十一届公共管理前沿论坛将于2026年10月11日在清华大学正式举行。本届前沿论坛主题围绕“数智时代的国家治理现代化”，并就人工智能、自主知识体系、新兴技术发展、“十五五”规划、碳双控等议题设置六大平行论坛进行探讨和交流。本届论坛邀请了来自全国各地的学术界与实务界代表参会，欢迎广大院友、师生参会！ 活动议程 上午：主论坛 签到 时 间： 10月11日（周日） 8:00-8:30 地 点： 清华大学蒙民伟音乐厅 开幕式（8:30-9:00） 主持人： 谢 矜 清华大学公共管理学院党委书记 致辞嘉宾： 杨 斌 清华大学校务委员会副主任，清华大学经济管理学院教授 路易斯·瓦熙（Luis Vassy） 巴黎政
\textbf{与你有关：} 这是当前用户可实际参与或申请的高可触达机会；参与范围：北京 / 线上；命中用户订阅版面：opportunity、ai、research；来源接近一手且权威度高；标题/正文包含高信号变化
\textbf{建议：} 检查报名资格、截止时间、参与形式和预期收益；合适则立即加入任务系统。
{\footnotesize\color{ormuted} 来源：北京智源社区活动 · \url{https://hub.baai.ac.cn/view/58571}}
\vspace{2mm}
\hrule height 0.2pt
\vspace{2mm}
{\large\bfseries Agentic AI / Generative AI}\par
Technical Blog Subscribe Agentic AI / Generative AI Oct 12, 2026 ReplaySSM: Rethinking State Storage for Faster Hybrid-Model Decoding Hybrid models interleave state-space model (SSM) layers with a few attention layers. NVIDIA Nemotron 3, Qwen3.5 through Qwen3.8, Kimi-K3, and GLM-5.3-Flash are... 14 MIN READ ReplaySSM: Rethinking State Storage for Faster Hybrid-Model Decoding Oct 08, 2026 Building Reliable Data Analyt
\textbf{与你有关：} 命中用户订阅版面：engineering、ai、business；来源接近一手且权威度高；标题/正文包含高信号变化
\textbf{建议：} 检查报名资格、截止时间、参与形式和预期收益；合适则立即加入任务系统。
{\footnotesize\color{ormuted} 来源：NVIDIA Developer Blog · \url{https://developer.nvidia.com/blog/category/generative-ai}}
\vspace{2mm}
\hrule height 0.2pt
\vspace{2mm}
{\large\bfseries One Model Family, Two Gold-Level Results: Fine-Tuning Nemotron for IOI and IMO}\par
Back to Articles One Model Family, Two Gold-Level Results: Fine-Tuning Nemotron for IOI and IMO Enterprise + Article Published October 7, 2026 Upvote 19 +13 Aleksander aficek Follow nvidia Igor Gitman igitman Follow nvidia Sean Narenthiran SeanNaren Follow nvidia Mehrzad Samadi mehrzads Follow nvidia Somshubra Majumdar smajumdar94 Follow nvidia The International Olympiad in Informatics (IOI) and the International Mat
\textbf{与你有关：} 命中用户订阅版面：engineering、ai、research；来源接近一手且权威度高；标题/正文包含高信号变化
\textbf{建议：} 检查报名资格、截止时间、参与形式和预期收益；合适则立即加入任务系统。
{\footnotesize\color{ormuted} 来源：Hugging Face Blog · \url{https://huggingface.co/blog/nvidia/nemotron-ioi-and-imo-2026}}
\vspace{2mm}
\hrule height 0.2pt
\vspace{2mm}
{\large\bfseries 报名中 2026-10-15 10:30:00 迈向通用化数据智能 LimiX-2}\par
报名中 2026-10-15 10:30:00 迈向通用化数据智能 LimiX-2
\textbf{与你有关：} 这是当前用户可实际参与或申请的高可触达机会；参与范围：北京 / 线上；命中用户订阅版面：opportunity、ai、research；来源接近一手且权威度高
\textbf{建议：} 检查报名资格、截止时间、参与形式和预期收益；合适则立即加入任务系统。
{\footnotesize\color{ormuted} 来源：北京智源社区活动 · \url{https://event.baai.ac.cn/activities/1088}}
\end{multicols}
\newpage
{\small\bfseries OR MORNING / SECTION 02 / 2026-10-13}
\vspace{2mm}
\hrule height 0.9pt
\vspace{1.5mm}
\hrule height 0.25pt
\vspace{4mm}
\section*{机会雷达}
{\color{ormuted} 北京、全国与线上真正可报名的 Hackathon、竞赛、科研招募、实习、Workshop 与截止。}
\vspace{2mm}
{\LARGE\bfseries 行政机构\_中国人民大学高瓴人工智能学院}\par
\vspace{1mm}
{\color{ormuted} 本版筛出 6 个高信号事件。}
\vspace{3mm}
{\Large\bfseries 行政机构\_中国人民大学高瓴人工智能学院}\par
首页 学院概况 学院简介 数说学院 现任领导 学术委员会 行政机构 新闻公告 学院新闻 学院公告 师生风采 光影视界 师资队伍 专任教师 博士后队伍 人才招聘 人才培养 本科生 学术型硕士研究生 专业学位硕士研究生 学术型博士研究生 专业学位博士研究生 在职课程培训班 同等学力申请硕士学位 科学研究 研究机构 研究成果 合作交流 对外合作 国际交流 校友名录 光影记录 校友名录 高礼研究院 EN 教师入口 学院概况 学院简介 数说学院 现任领导 学术委员会 行政机构 行政机构 您当前所在的位置： 首页 > 学院概况 > 行政机构 教务工作 - 教务管理：负责排课、选课、成绩管理，以及毕业实践和毕业论文的组织与管理，包括出具成绩单和协助学历研究生招生工作。 - 教师教学工作管理：计算和考核教师的教学工作量，确保教学任务的合理分配和履行。 - 教学委员会支持：协助各教学工作委员会的秘书处工作，包括会议组织、文件管理等。 - 学科建设
\textbf{与你有关：} 这是当前用户可实际参与或申请的高可触达机会；参与范围：北京 / 人大；命中用户订阅版面：opportunity、research、ai；来源接近一手且权威度高；标题/正文包含高信号变化
\textbf{建议：} 检查报名资格、截止时间、参与形式和预期收益；合适则立即加入任务系统。
{\footnotesize\color{ormuted} 来源：中国人民大学高瓴人工智能学院公告 · \url{https://ai.ruc.edu.cn/xygk/xzjg/index.htm}}
\vspace{4mm}
\hrule height 0.25pt
\vspace{3mm}
\begin{multicols}{2}
{\large\bfseries 姚班培养-清华大学交叉信息研究院}\par
首页 热点新闻 讲座预告 院系概况 院系简介 院长寄语 现任领导 院系沿革 新闻动态 院系动态 媒体聚焦 人员队伍 全职教师 研究系列 荣誉客座 兼职教师 博士后 行政人员 实验员 姚班本科生 姚班培养 姚班招生 研究生 研究生培养 研究生招生 量子中心 中心简介 师资力量 科研概况 量子中心动态 科学研究 组织架构 课题组介绍 研究动态 科研简报 校园生活 学生活动 衣食住行 公开招聘 教师招聘 博士后招聘 职员招聘 校友频道 校友风采 校友俱乐部 校友捐赠 姚班本科生 姚班培养 姚班招生 清华学堂计算机科学实验班 姚班介绍 “清华学堂计算机科学实验班（姚班）”由世界著名计算机科学家姚期智院士于2005年创办，是清华大学培养国际顶尖计算机科学及交叉创新人才的本科生实验班，被誉为“最优秀的本科生和最优秀的本科教育”标杆。 姚班面向科技发展前沿与国家重大需求构建培养体系，以计算机科学、人工智能和量子信息三大前沿专业方向为核心，供学
\textbf{与你有关：} 这是当前用户可实际参与或申请的高可触达机会；参与范围：北京 / 清华；命中用户订阅版面：opportunity、research、ai；来源接近一手且权威度高；标题/正文包含高信号变化
\textbf{建议：} 检查报名资格、截止时间、参与形式和预期收益；合适则立即加入任务系统。
{\footnotesize\color{ormuted} 来源：清华交叉信息研究院 · \url{https://iiis.tsinghua.edu.cn/ybbks/ybpy.htm}}
\vspace{2mm}
\hrule height 0.2pt
\vspace{2mm}
{\large\bfseries 中心简介-清华大学交叉信息研究院}\par
首页 热点新闻 讲座预告 院系概况 院系简介 院长寄语 现任领导 院系沿革 新闻动态 院系动态 媒体聚焦 人员队伍 全职教师 研究系列 荣誉客座 兼职教师 博士后 行政人员 实验员 姚班本科生 姚班培养 姚班招生 研究生 研究生培养 研究生招生 量子中心 中心简介 师资力量 科研概况 量子中心动态 科学研究 组织架构 课题组介绍 研究动态 科研简报 校园生活 学生活动 衣食住行 公开招聘 教师招聘 博士后招聘 职员招聘 校友频道 校友风采 校友俱乐部 校友捐赠 量子中心 中心简介 师资力量 科研概况 量子中心动态 量子中心 中心介绍 量子信息是一门新兴的交叉学科领域，涉及计算机科学、信息科学和量子物理，为后摩尔时代提供了新的技术路径，探索了21世纪信息科学发展的新基础与新方法。量子信息技术已成为各国战略竞争的焦点之一。量子信息科学作为量子调控的核心，是信息网络与纳米科学等领域的交叉方向，代表了中国未来科技发展的重要战略方向之一。
\textbf{与你有关：} 这是当前用户可实际参与或申请的高可触达机会；参与范围：北京 / 清华；命中用户订阅版面：opportunity、research、ai；来源接近一手且权威度高；标题/正文包含高信号变化
\textbf{建议：} 检查报名资格、截止时间、参与形式和预期收益；合适则立即加入任务系统。
{\footnotesize\color{ormuted} 来源：清华交叉信息研究院 · \url{https://iiis.tsinghua.edu.cn/lzzx/zxjj.htm}}
\vspace{2mm}
\hrule height 0.2pt
\vspace{2mm}
{\large\bfseries 最新日程：清华大学第二十一届公共管理前沿论坛即将开幕}\par
最新日程：清华大学第二十一届公共管理前沿论坛即将开幕 IR AI SoftEng 清华大学人工智能国际治理研究院 2026-10-10 22:50 分享 以下文章来源于mp.weixin.qq.com 点击蓝字 关注我们 清华大学第二十一届公共管理前沿论坛将于2026年10月11日在清华大学正式举行。本届前沿论坛主题围绕“数智时代的国家治理现代化”，并就人工智能、自主知识体系、新兴技术发展、“十五五”规划、碳双控等议题设置六大平行论坛进行探讨和交流。本届论坛邀请了来自全国各地的学术界与实务界代表参会，欢迎广大院友、师生参会！ 活动议程 上午：主论坛 签到 时 间： 10月11日（周日） 8:00-8:30 地 点： 清华大学蒙民伟音乐厅 开幕式（8:30-9:00） 主持人： 谢 矜 清华大学公共管理学院党委书记 致辞嘉宾： 杨 斌 清华大学校务委员会副主任，清华大学经济管理学院教授 路易斯·瓦熙（Luis Vassy） 巴黎政
\textbf{与你有关：} 这是当前用户可实际参与或申请的高可触达机会；参与范围：北京 / 线上；命中用户订阅版面：opportunity、ai、research；来源接近一手且权威度高；标题/正文包含高信号变化
\textbf{建议：} 检查报名资格、截止时间、参与形式和预期收益；合适则立即加入任务系统。
{\footnotesize\color{ormuted} 来源：北京智源社区活动 · \url{https://hub.baai.ac.cn/view/58571}}
\vspace{2mm}
\hrule height 0.2pt
\vspace{2mm}
{\large\bfseries 报名中 2026-10-15 10:30:00 迈向通用化数据智能 LimiX-2}\par
报名中 2026-10-15 10:30:00 迈向通用化数据智能 LimiX-2
\textbf{与你有关：} 这是当前用户可实际参与或申请的高可触达机会；参与范围：北京 / 线上；命中用户订阅版面：opportunity、ai、research；来源接近一手且权威度高
\textbf{建议：} 检查报名资格、截止时间、参与形式和预期收益；合适则立即加入任务系统。
{\footnotesize\color{ormuted} 来源：北京智源社区活动 · \url{https://event.baai.ac.cn/activities/1088}}
\vspace{2mm}
\hrule height 0.2pt
\vspace{2mm}
{\large\bfseries Select Year: (2026)}\par
Select Year: (2026) 2026 2025 2024 2023 2022 2021 2020 2019 2018 2017 2016 2015 2014 2013 2012 2011 2010 2009 2008 2007 2006 Earlier Conferences
\textbf{与你有关：} 命中用户订阅版面：research、opportunity、ai；来源接近一手且权威度高；标题/正文包含高信号变化
\textbf{建议：} 检查报名资格、截止时间、参与形式和预期收益；合适则立即加入任务系统。
{\footnotesize\color{ormuted} 来源：NeurIPS · \url{https://neurips.cc/}}
\end{multicols}
\newpage
{\small\bfseries OR MORNING / SECTION 03 / 2026-10-13}
\vspace{2mm}
\hrule height 0.9pt
\vspace{1.5mm}
\hrule height 0.25pt
\vspace{4mm}
\section*{科研前沿}
{\color{ormuted} 真正会改变你接下来怎么做研究、怎么写、怎么投的信号。}
\vspace{2mm}
{\LARGE\bfseries 行政机构\_中国人民大学高瓴人工智能学院}\par
\vspace{1mm}
{\color{ormuted} 本版筛出 5 个高信号事件。}
\vspace{3mm}
{\Large\bfseries 行政机构\_中国人民大学高瓴人工智能学院}\par
首页 学院概况 学院简介 数说学院 现任领导 学术委员会 行政机构 新闻公告 学院新闻 学院公告 师生风采 光影视界 师资队伍 专任教师 博士后队伍 人才招聘 人才培养 本科生 学术型硕士研究生 专业学位硕士研究生 学术型博士研究生 专业学位博士研究生 在职课程培训班 同等学力申请硕士学位 科学研究 研究机构 研究成果 合作交流 对外合作 国际交流 校友名录 光影记录 校友名录 高礼研究院 EN 教师入口 学院概况 学院简介 数说学院 现任领导 学术委员会 行政机构 行政机构 您当前所在的位置： 首页 > 学院概况 > 行政机构 教务工作 - 教务管理：负责排课、选课、成绩管理，以及毕业实践和毕业论文的组织与管理，包括出具成绩单和协助学历研究生招生工作。 - 教师教学工作管理：计算和考核教师的教学工作量，确保教学任务的合理分配和履行。 - 教学委员会支持：协助各教学工作委员会的秘书处工作，包括会议组织、文件管理等。 - 学科建设
\textbf{与你有关：} 这是当前用户可实际参与或申请的高可触达机会；参与范围：北京 / 人大；命中用户订阅版面：opportunity、research、ai；来源接近一手且权威度高；标题/正文包含高信号变化
\textbf{建议：} 检查报名资格、截止时间、参与形式和预期收益；合适则立即加入任务系统。
{\footnotesize\color{ormuted} 来源：中国人民大学高瓴人工智能学院公告 · \url{https://ai.ruc.edu.cn/xygk/xzjg/index.htm}}
\vspace{4mm}
\hrule height 0.25pt
\vspace{3mm}
\begin{multicols}{2}
{\large\bfseries 姚班培养-清华大学交叉信息研究院}\par
首页 热点新闻 讲座预告 院系概况 院系简介 院长寄语 现任领导 院系沿革 新闻动态 院系动态 媒体聚焦 人员队伍 全职教师 研究系列 荣誉客座 兼职教师 博士后 行政人员 实验员 姚班本科生 姚班培养 姚班招生 研究生 研究生培养 研究生招生 量子中心 中心简介 师资力量 科研概况 量子中心动态 科学研究 组织架构 课题组介绍 研究动态 科研简报 校园生活 学生活动 衣食住行 公开招聘 教师招聘 博士后招聘 职员招聘 校友频道 校友风采 校友俱乐部 校友捐赠 姚班本科生 姚班培养 姚班招生 清华学堂计算机科学实验班 姚班介绍 “清华学堂计算机科学实验班（姚班）”由世界著名计算机科学家姚期智院士于2005年创办，是清华大学培养国际顶尖计算机科学及交叉创新人才的本科生实验班，被誉为“最优秀的本科生和最优秀的本科教育”标杆。 姚班面向科技发展前沿与国家重大需求构建培养体系，以计算机科学、人工智能和量子信息三大前沿专业方向为核心，供学
\textbf{与你有关：} 这是当前用户可实际参与或申请的高可触达机会；参与范围：北京 / 清华；命中用户订阅版面：opportunity、research、ai；来源接近一手且权威度高；标题/正文包含高信号变化
\textbf{建议：} 检查报名资格、截止时间、参与形式和预期收益；合适则立即加入任务系统。
{\footnotesize\color{ormuted} 来源：清华交叉信息研究院 · \url{https://iiis.tsinghua.edu.cn/ybbks/ybpy.htm}}
\vspace{2mm}
\hrule height 0.2pt
\vspace{2mm}
{\large\bfseries 中心简介-清华大学交叉信息研究院}\par
首页 热点新闻 讲座预告 院系概况 院系简介 院长寄语 现任领导 院系沿革 新闻动态 院系动态 媒体聚焦 人员队伍 全职教师 研究系列 荣誉客座 兼职教师 博士后 行政人员 实验员 姚班本科生 姚班培养 姚班招生 研究生 研究生培养 研究生招生 量子中心 中心简介 师资力量 科研概况 量子中心动态 科学研究 组织架构 课题组介绍 研究动态 科研简报 校园生活 学生活动 衣食住行 公开招聘 教师招聘 博士后招聘 职员招聘 校友频道 校友风采 校友俱乐部 校友捐赠 量子中心 中心简介 师资力量 科研概况 量子中心动态 量子中心 中心介绍 量子信息是一门新兴的交叉学科领域，涉及计算机科学、信息科学和量子物理，为后摩尔时代提供了新的技术路径，探索了21世纪信息科学发展的新基础与新方法。量子信息技术已成为各国战略竞争的焦点之一。量子信息科学作为量子调控的核心，是信息网络与纳米科学等领域的交叉方向，代表了中国未来科技发展的重要战略方向之一。
\textbf{与你有关：} 这是当前用户可实际参与或申请的高可触达机会；参与范围：北京 / 清华；命中用户订阅版面：opportunity、research、ai；来源接近一手且权威度高；标题/正文包含高信号变化
\textbf{建议：} 检查报名资格、截止时间、参与形式和预期收益；合适则立即加入任务系统。
{\footnotesize\color{ormuted} 来源：清华交叉信息研究院 · \url{https://iiis.tsinghua.edu.cn/lzzx/zxjj.htm}}
\vspace{2mm}
\hrule height 0.2pt
\vspace{2mm}
{\large\bfseries 最新日程：清华大学第二十一届公共管理前沿论坛即将开幕}\par
最新日程：清华大学第二十一届公共管理前沿论坛即将开幕 IR AI SoftEng 清华大学人工智能国际治理研究院 2026-10-10 22:50 分享 以下文章来源于mp.weixin.qq.com 点击蓝字 关注我们 清华大学第二十一届公共管理前沿论坛将于2026年10月11日在清华大学正式举行。本届前沿论坛主题围绕“数智时代的国家治理现代化”，并就人工智能、自主知识体系、新兴技术发展、“十五五”规划、碳双控等议题设置六大平行论坛进行探讨和交流。本届论坛邀请了来自全国各地的学术界与实务界代表参会，欢迎广大院友、师生参会！ 活动议程 上午：主论坛 签到 时 间： 10月11日（周日） 8:00-8:30 地 点： 清华大学蒙民伟音乐厅 开幕式（8:30-9:00） 主持人： 谢 矜 清华大学公共管理学院党委书记 致辞嘉宾： 杨 斌 清华大学校务委员会副主任，清华大学经济管理学院教授 路易斯·瓦熙（Luis Vassy） 巴黎政
\textbf{与你有关：} 这是当前用户可实际参与或申请的高可触达机会；参与范围：北京 / 线上；命中用户订阅版面：opportunity、ai、research；来源接近一手且权威度高；标题/正文包含高信号变化
\textbf{建议：} 检查报名资格、截止时间、参与形式和预期收益；合适则立即加入任务系统。
{\footnotesize\color{ormuted} 来源：北京智源社区活动 · \url{https://hub.baai.ac.cn/view/58571}}
\vspace{2mm}
\hrule height 0.2pt
\vspace{2mm}
{\large\bfseries 报名中 2026-10-15 10:30:00 迈向通用化数据智能 LimiX-2}\par
报名中 2026-10-15 10:30:00 迈向通用化数据智能 LimiX-2
\textbf{与你有关：} 这是当前用户可实际参与或申请的高可触达机会；参与范围：北京 / 线上；命中用户订阅版面：opportunity、ai、research；来源接近一手且权威度高
\textbf{建议：} 检查报名资格、截止时间、参与形式和预期收益；合适则立即加入任务系统。
{\footnotesize\color{ormuted} 来源：北京智源社区活动 · \url{https://event.baai.ac.cn/activities/1088}}
\end{multicols}
\newpage
{\small\bfseries OR MORNING / SECTION 04 / 2026-10-13}
\vspace{2mm}
\hrule height 0.9pt
\vspace{1.5mm}
\hrule height 0.25pt
\vspace{4mm}
\section*{AI Frontier}
{\color{ormuted} Frontier model、Agent、RL、Reasoning、Safety。}
\vspace{2mm}
{\LARGE\bfseries 行政机构\_中国人民大学高瓴人工智能学院}\par
\vspace{1mm}
{\color{ormuted} 本版筛出 5 个高信号事件。}
\vspace{3mm}
{\Large\bfseries 行政机构\_中国人民大学高瓴人工智能学院}\par
首页 学院概况 学院简介 数说学院 现任领导 学术委员会 行政机构 新闻公告 学院新闻 学院公告 师生风采 光影视界 师资队伍 专任教师 博士后队伍 人才招聘 人才培养 本科生 学术型硕士研究生 专业学位硕士研究生 学术型博士研究生 专业学位博士研究生 在职课程培训班 同等学力申请硕士学位 科学研究 研究机构 研究成果 合作交流 对外合作 国际交流 校友名录 光影记录 校友名录 高礼研究院 EN 教师入口 学院概况 学院简介 数说学院 现任领导 学术委员会 行政机构 行政机构 您当前所在的位置： 首页 > 学院概况 > 行政机构 教务工作 - 教务管理：负责排课、选课、成绩管理，以及毕业实践和毕业论文的组织与管理，包括出具成绩单和协助学历研究生招生工作。 - 教师教学工作管理：计算和考核教师的教学工作量，确保教学任务的合理分配和履行。 - 教学委员会支持：协助各教学工作委员会的秘书处工作，包括会议组织、文件管理等。 - 学科建设
\textbf{与你有关：} 这是当前用户可实际参与或申请的高可触达机会；参与范围：北京 / 人大；命中用户订阅版面：opportunity、research、ai；来源接近一手且权威度高；标题/正文包含高信号变化
\textbf{建议：} 检查报名资格、截止时间、参与形式和预期收益；合适则立即加入任务系统。
{\footnotesize\color{ormuted} 来源：中国人民大学高瓴人工智能学院公告 · \url{https://ai.ruc.edu.cn/xygk/xzjg/index.htm}}
\vspace{4mm}
\hrule height 0.25pt
\vspace{3mm}
\begin{multicols}{2}
{\large\bfseries 姚班培养-清华大学交叉信息研究院}\par
首页 热点新闻 讲座预告 院系概况 院系简介 院长寄语 现任领导 院系沿革 新闻动态 院系动态 媒体聚焦 人员队伍 全职教师 研究系列 荣誉客座 兼职教师 博士后 行政人员 实验员 姚班本科生 姚班培养 姚班招生 研究生 研究生培养 研究生招生 量子中心 中心简介 师资力量 科研概况 量子中心动态 科学研究 组织架构 课题组介绍 研究动态 科研简报 校园生活 学生活动 衣食住行 公开招聘 教师招聘 博士后招聘 职员招聘 校友频道 校友风采 校友俱乐部 校友捐赠 姚班本科生 姚班培养 姚班招生 清华学堂计算机科学实验班 姚班介绍 “清华学堂计算机科学实验班（姚班）”由世界著名计算机科学家姚期智院士于2005年创办，是清华大学培养国际顶尖计算机科学及交叉创新人才的本科生实验班，被誉为“最优秀的本科生和最优秀的本科教育”标杆。 姚班面向科技发展前沿与国家重大需求构建培养体系，以计算机科学、人工智能和量子信息三大前沿专业方向为核心，供学
\textbf{与你有关：} 这是当前用户可实际参与或申请的高可触达机会；参与范围：北京 / 清华；命中用户订阅版面：opportunity、research、ai；来源接近一手且权威度高；标题/正文包含高信号变化
\textbf{建议：} 检查报名资格、截止时间、参与形式和预期收益；合适则立即加入任务系统。
{\footnotesize\color{ormuted} 来源：清华交叉信息研究院 · \url{https://iiis.tsinghua.edu.cn/ybbks/ybpy.htm}}
\vspace{2mm}
\hrule height 0.2pt
\vspace{2mm}
{\large\bfseries 中心简介-清华大学交叉信息研究院}\par
首页 热点新闻 讲座预告 院系概况 院系简介 院长寄语 现任领导 院系沿革 新闻动态 院系动态 媒体聚焦 人员队伍 全职教师 研究系列 荣誉客座 兼职教师 博士后 行政人员 实验员 姚班本科生 姚班培养 姚班招生 研究生 研究生培养 研究生招生 量子中心 中心简介 师资力量 科研概况 量子中心动态 科学研究 组织架构 课题组介绍 研究动态 科研简报 校园生活 学生活动 衣食住行 公开招聘 教师招聘 博士后招聘 职员招聘 校友频道 校友风采 校友俱乐部 校友捐赠 量子中心 中心简介 师资力量 科研概况 量子中心动态 量子中心 中心介绍 量子信息是一门新兴的交叉学科领域，涉及计算机科学、信息科学和量子物理，为后摩尔时代提供了新的技术路径，探索了21世纪信息科学发展的新基础与新方法。量子信息技术已成为各国战略竞争的焦点之一。量子信息科学作为量子调控的核心，是信息网络与纳米科学等领域的交叉方向，代表了中国未来科技发展的重要战略方向之一。
\textbf{与你有关：} 这是当前用户可实际参与或申请的高可触达机会；参与范围：北京 / 清华；命中用户订阅版面：opportunity、research、ai；来源接近一手且权威度高；标题/正文包含高信号变化
\textbf{建议：} 检查报名资格、截止时间、参与形式和预期收益；合适则立即加入任务系统。
{\footnotesize\color{ormuted} 来源：清华交叉信息研究院 · \url{https://iiis.tsinghua.edu.cn/lzzx/zxjj.htm}}
\vspace{2mm}
\hrule height 0.2pt
\vspace{2mm}
{\large\bfseries 最新日程：清华大学第二十一届公共管理前沿论坛即将开幕}\par
最新日程：清华大学第二十一届公共管理前沿论坛即将开幕 IR AI SoftEng 清华大学人工智能国际治理研究院 2026-10-10 22:50 分享 以下文章来源于mp.weixin.qq.com 点击蓝字 关注我们 清华大学第二十一届公共管理前沿论坛将于2026年10月11日在清华大学正式举行。本届前沿论坛主题围绕“数智时代的国家治理现代化”，并就人工智能、自主知识体系、新兴技术发展、“十五五”规划、碳双控等议题设置六大平行论坛进行探讨和交流。本届论坛邀请了来自全国各地的学术界与实务界代表参会，欢迎广大院友、师生参会！ 活动议程 上午：主论坛 签到 时 间： 10月11日（周日） 8:00-8:30 地 点： 清华大学蒙民伟音乐厅 开幕式（8:30-9:00） 主持人： 谢 矜 清华大学公共管理学院党委书记 致辞嘉宾： 杨 斌 清华大学校务委员会副主任，清华大学经济管理学院教授 路易斯·瓦熙（Luis Vassy） 巴黎政
\textbf{与你有关：} 这是当前用户可实际参与或申请的高可触达机会；参与范围：北京 / 线上；命中用户订阅版面：opportunity、ai、research；来源接近一手且权威度高；标题/正文包含高信号变化
\textbf{建议：} 检查报名资格、截止时间、参与形式和预期收益；合适则立即加入任务系统。
{\footnotesize\color{ormuted} 来源：北京智源社区活动 · \url{https://hub.baai.ac.cn/view/58571}}
\vspace{2mm}
\hrule height 0.2pt
\vspace{2mm}
{\large\bfseries 报名中 2026-10-15 10:30:00 迈向通用化数据智能 LimiX-2}\par
报名中 2026-10-15 10:30:00 迈向通用化数据智能 LimiX-2
\textbf{与你有关：} 这是当前用户可实际参与或申请的高可触达机会；参与范围：北京 / 线上；命中用户订阅版面：opportunity、ai、research；来源接近一手且权威度高
\textbf{建议：} 检查报名资格、截止时间、参与形式和预期收益；合适则立即加入任务系统。
{\footnotesize\color{ormuted} 来源：北京智源社区活动 · \url{https://event.baai.ac.cn/activities/1088}}
\end{multicols}
\newpage
{\small\bfseries OR MORNING / SECTION 05 / 2026-10-13}
\vspace{2mm}
\hrule height 0.9pt
\vspace{1.5mm}
\hrule height 0.25pt
\vspace{4mm}
\section*{数学 × AI}
{\color{ormuted} AI for Mathematics、Lean、证明与重要数学进展。}
\vspace{2mm}
{\LARGE\bfseries One Model Family, Two Gold-Level Results: Fine-Tuning Nemotron for IOI and IMO}\par
\vspace{1mm}
{\color{ormuted} 本版筛出 5 个高信号事件。}
\vspace{3mm}
{\Large\bfseries One Model Family, Two Gold-Level Results: Fine-Tuning Nemotron for IOI and IMO}\par
Back to Articles One Model Family, Two Gold-Level Results: Fine-Tuning Nemotron for IOI and IMO Enterprise + Article Published October 7, 2026 Upvote 19 +13 Aleksander aficek Follow nvidia Igor Gitman igitman Follow nvidia Sean Narenthiran SeanNaren Follow nvidia Mehrzad Samadi mehrzads Follow nvidia Somshubra Majumdar smajumdar94 Follow nvidia The International Olympiad in Informatics (IOI) and the International Mat
\textbf{与你有关：} 命中用户订阅版面：engineering、ai、research；来源接近一手且权威度高；标题/正文包含高信号变化
\textbf{建议：} 检查报名资格、截止时间、参与形式和预期收益；合适则立即加入任务系统。
{\footnotesize\color{ormuted} 来源：Hugging Face Blog · \url{https://huggingface.co/blog/nvidia/nemotron-ioi-and-imo-2026}}
\vspace{4mm}
\hrule height 0.25pt
\vspace{3mm}
\begin{multicols}{2}
{\large\bfseries Beijing Institute of Mathematical Sciences and Applications}\par
31st August, 2026 The Review of Mathematical Economics vol.1 (2026), no.2 has been published 28th August, 2026 Congratulations! Prof. Yaqing Wang Selected for IJCAI 2026 Early Career Spotlights 29th July, 2026 "Feiyan" Makes Successful Maiden Flight 25th June, 2026 When Materials Science Meets Mathematics: How to Understand Complex Materials 25th June, 2026 An Exclusive Interview with Fellow Oluwole Daniel Makinde 22
\textbf{与你有关：} 标题/正文包含高信号变化
\textbf{建议：} 检查报名资格、截止时间、参与形式和预期收益；合适则立即加入任务系统。
{\footnotesize\color{ormuted} 来源：BIMSA · \url{https://www.bimsa.cn/}}
\vspace{2mm}
\hrule height 0.2pt
\vspace{2mm}
{\large\bfseries Skip to main content}\par
Research Our research teams investigate the safety, inner workings, and societal impacts of AI models—so that artificial intelligence has a positive impact as it becomes increasingly capable. Research teams: Alignment Economics Interpretability Societal Impacts Frontier Red Team Alignment The Alignment team works to understand the risks of AI models and develop ways to ensure that future ones remain helpful, honest, 
\textbf{与你有关：} 命中用户订阅版面：ai、research、math；来源接近一手且权威度高；标题/正文包含高信号变化
{\footnotesize\color{ormuted} 来源：Anthropic Research · \url{https://www.anthropic.com/research}}
\vspace{2mm}
\hrule height 0.2pt
\vspace{2mm}
{\large\bfseries Societal Impacts Sep 29, 2026 What do you want from AI? We’re launching a new study using Anthropic Interviewer to learn from your experiences with AI.}\par
Societal Impacts What do you want from AI? Sep 29, 2026 Participate here We’re launching a new study using Anthropic Interviewer to learn from your experiences with AI, and we’d like you to participate. After you finish, you can decide to make your interview public, so that anyone, not just Anthropic, can read and learn from it. You can participate here . We are at a pivotal moment in the development of AI, as its gr
\textbf{与你有关：} 命中用户订阅版面：ai、research、math；来源接近一手且权威度高；标题/正文包含高信号变化
{\footnotesize\color{ormuted} 来源：Anthropic Research · \url{https://www.anthropic.com/research/your-thoughts-on-ai}}
\vspace{2mm}
\hrule height 0.2pt
\vspace{2mm}
{\large\bfseries Oct 9, 2026 Alignment Investigating unintended model actions in our evaluations and internal use}\par
Oct 9, 2026 Alignment Investigating unintended model actions in our evaluations and internal use
\textbf{与你有关：} 命中用户订阅版面：ai、research、math；来源接近一手且权威度高；标题/正文包含高信号变化
{\footnotesize\color{ormuted} 来源：Anthropic Research · \url{https://www.anthropic.com/research/investigating-unintended-model-actions}}
\end{multicols}
\newpage
{\small\bfseries OR MORNING / SECTION 06 / 2026-10-13}
\vspace{2mm}
\hrule height 0.9pt
\vspace{1.5mm}
\hrule height 0.25pt
\vspace{4mm}
\section*{Engineering}
{\color{ormuted} 开源、框架、GPU、推理基础设施与真正值得使用的新工具。}
\vspace{2mm}
{\LARGE\bfseries 最新日程：清华大学第二十一届公共管理前沿论坛即将开幕}\par
\vspace{1mm}
{\color{ormuted} 本版筛出 5 个高信号事件。}
\vspace{3mm}
{\Large\bfseries 最新日程：清华大学第二十一届公共管理前沿论坛即将开幕}\par
最新日程：清华大学第二十一届公共管理前沿论坛即将开幕 IR AI SoftEng 清华大学人工智能国际治理研究院 2026-10-10 22:50 分享 以下文章来源于mp.weixin.qq.com 点击蓝字 关注我们 清华大学第二十一届公共管理前沿论坛将于2026年10月11日在清华大学正式举行。本届前沿论坛主题围绕“数智时代的国家治理现代化”，并就人工智能、自主知识体系、新兴技术发展、“十五五”规划、碳双控等议题设置六大平行论坛进行探讨和交流。本届论坛邀请了来自全国各地的学术界与实务界代表参会，欢迎广大院友、师生参会！ 活动议程 上午：主论坛 签到 时 间： 10月11日（周日） 8:00-8:30 地 点： 清华大学蒙民伟音乐厅 开幕式（8:30-9:00） 主持人： 谢 矜 清华大学公共管理学院党委书记 致辞嘉宾： 杨 斌 清华大学校务委员会副主任，清华大学经济管理学院教授 路易斯·瓦熙（Luis Vassy） 巴黎政
\textbf{与你有关：} 这是当前用户可实际参与或申请的高可触达机会；参与范围：北京 / 线上；命中用户订阅版面：opportunity、ai、research；来源接近一手且权威度高；标题/正文包含高信号变化
\textbf{建议：} 检查报名资格、截止时间、参与形式和预期收益；合适则立即加入任务系统。
{\footnotesize\color{ormuted} 来源：北京智源社区活动 · \url{https://hub.baai.ac.cn/view/58571}}
\vspace{4mm}
\hrule height 0.25pt
\vspace{3mm}
\begin{multicols}{2}
{\large\bfseries 报名中 2026-10-15 10:30:00 迈向通用化数据智能 LimiX-2}\par
报名中 2026-10-15 10:30:00 迈向通用化数据智能 LimiX-2
\textbf{与你有关：} 这是当前用户可实际参与或申请的高可触达机会；参与范围：北京 / 线上；命中用户订阅版面：opportunity、ai、research；来源接近一手且权威度高
\textbf{建议：} 检查报名资格、截止时间、参与形式和预期收益；合适则立即加入任务系统。
{\footnotesize\color{ormuted} 来源：北京智源社区活动 · \url{https://event.baai.ac.cn/activities/1088}}
\vspace{2mm}
\hrule height 0.2pt
\vspace{2mm}
{\large\bfseries Agentic AI / Generative AI}\par
Technical Blog Subscribe Agentic AI / Generative AI Oct 12, 2026 ReplaySSM: Rethinking State Storage for Faster Hybrid-Model Decoding Hybrid models interleave state-space model (SSM) layers with a few attention layers. NVIDIA Nemotron 3, Qwen3.5 through Qwen3.8, Kimi-K3, and GLM-5.3-Flash are... 14 MIN READ ReplaySSM: Rethinking State Storage for Faster Hybrid-Model Decoding Oct 08, 2026 Building Reliable Data Analyt
\textbf{与你有关：} 命中用户订阅版面：engineering、ai、business；来源接近一手且权威度高；标题/正文包含高信号变化
\textbf{建议：} 检查报名资格、截止时间、参与形式和预期收益；合适则立即加入任务系统。
{\footnotesize\color{ormuted} 来源：NVIDIA Developer Blog · \url{https://developer.nvidia.com/blog/category/generative-ai}}
\vspace{2mm}
\hrule height 0.2pt
\vspace{2mm}
{\large\bfseries One Model Family, Two Gold-Level Results: Fine-Tuning Nemotron for IOI and IMO}\par
Back to Articles One Model Family, Two Gold-Level Results: Fine-Tuning Nemotron for IOI and IMO Enterprise + Article Published October 7, 2026 Upvote 19 +13 Aleksander aficek Follow nvidia Igor Gitman igitman Follow nvidia Sean Narenthiran SeanNaren Follow nvidia Mehrzad Samadi mehrzads Follow nvidia Somshubra Majumdar smajumdar94 Follow nvidia The International Olympiad in Informatics (IOI) and the International Mat
\textbf{与你有关：} 命中用户订阅版面：engineering、ai、research；来源接近一手且权威度高；标题/正文包含高信号变化
\textbf{建议：} 检查报名资格、截止时间、参与形式和预期收益；合适则立即加入任务系统。
{\footnotesize\color{ormuted} 来源：Hugging Face Blog · \url{https://huggingface.co/blog/nvidia/nemotron-ioi-and-imo-2026}}
\vspace{2mm}
\hrule height 0.2pt
\vspace{2mm}
{\large\bfseries Back to Careers}\par
Back to Careers Join our team All Teams AI Research \& Engineering Applied AI Communications Compute Data Science \& Analytics Engineering \& Design - Product Finance Hardware Legal Marketing \& Brand People Product Management, Support, \& Operations Public Benefit Public Policy Safeguards (Trust \& Safety) Sales Security Software Engineering - Infrastructure Technical Education Technical Program Management All Offices Ban
\textbf{与你有关：} 参与范围：全球 / 多为海外全职；命中用户订阅版面：opportunity、ai、research；来源接近一手且权威度高；标题/正文包含高信号变化
\textbf{建议：} 检查报名资格、截止时间、参与形式和预期收益；合适则立即加入任务系统。
{\footnotesize\color{ormuted} 来源：Anthropic Careers · \url{https://www.anthropic.com/careers}}
\end{multicols}
\newpage
{\small\bfseries OR MORNING / SECTION 07 / 2026-10-13}
\vspace{2mm}
\hrule height 0.9pt
\vspace{1.5mm}
\hrule height 0.25pt
\vspace{4mm}
\section*{AI 商业 / 产业}
{\color{ormuted} 公司、融资、市场与可能改变 AI 发展路径的产业信号。}
\vspace{2mm}
{\LARGE\bfseries Agentic AI / Generative AI}\par
\vspace{1mm}
{\color{ormuted} 本版筛出 5 个高信号事件。}
\vspace{3mm}
{\Large\bfseries Agentic AI / Generative AI}\par
Technical Blog Subscribe Agentic AI / Generative AI Oct 12, 2026 ReplaySSM: Rethinking State Storage for Faster Hybrid-Model Decoding Hybrid models interleave state-space model (SSM) layers with a few attention layers. NVIDIA Nemotron 3, Qwen3.5 through Qwen3.8, Kimi-K3, and GLM-5.3-Flash are... 14 MIN READ ReplaySSM: Rethinking State Storage for Faster Hybrid-Model Decoding Oct 08, 2026 Building Reliable Data Analyt
\textbf{与你有关：} 命中用户订阅版面：engineering、ai、business；来源接近一手且权威度高；标题/正文包含高信号变化
\textbf{建议：} 检查报名资格、截止时间、参与形式和预期收益；合适则立即加入任务系统。
{\footnotesize\color{ormuted} 来源：NVIDIA Developer Blog · \url{https://developer.nvidia.com/blog/category/generative-ai}}
\vspace{4mm}
\hrule height 0.25pt
\vspace{3mm}
\begin{multicols}{2}
{\large\bfseries One Model Family, Two Gold-Level Results: Fine-Tuning Nemotron for IOI and IMO}\par
Back to Articles One Model Family, Two Gold-Level Results: Fine-Tuning Nemotron for IOI and IMO Enterprise + Article Published October 7, 2026 Upvote 19 +13 Aleksander aficek Follow nvidia Igor Gitman igitman Follow nvidia Sean Narenthiran SeanNaren Follow nvidia Mehrzad Samadi mehrzads Follow nvidia Somshubra Majumdar smajumdar94 Follow nvidia The International Olympiad in Informatics (IOI) and the International Mat
\textbf{与你有关：} 命中用户订阅版面：engineering、ai、research；来源接近一手且权威度高；标题/正文包含高信号变化
\textbf{建议：} 检查报名资格、截止时间、参与形式和预期收益；合适则立即加入任务系统。
{\footnotesize\color{ormuted} 来源：Hugging Face Blog · \url{https://huggingface.co/blog/nvidia/nemotron-ioi-and-imo-2026}}
\vspace{2mm}
\hrule height 0.2pt
\vspace{2mm}
{\large\bfseries Back to Careers}\par
Back to Careers Join our team All Teams AI Research \& Engineering Applied AI Communications Compute Data Science \& Analytics Engineering \& Design - Product Finance Hardware Legal Marketing \& Brand People Product Management, Support, \& Operations Public Benefit Public Policy Safeguards (Trust \& Safety) Sales Security Software Engineering - Infrastructure Technical Education Technical Program Management All Offices Ban
\textbf{与你有关：} 参与范围：全球 / 多为海外全职；命中用户订阅版面：opportunity、ai、research；来源接近一手且权威度高；标题/正文包含高信号变化
\textbf{建议：} 检查报名资格、截止时间、参与形式和预期收益；合适则立即加入任务系统。
{\footnotesize\color{ormuted} 来源：Anthropic Careers · \url{https://www.anthropic.com/careers}}
\vspace{2mm}
\hrule height 0.2pt
\vspace{2mm}
{\large\bfseries EDIC for Agri-Food launched to support digital innovation across Europe’s food and agriculture sector}\par
EDIC for Agri-Food launched to support digital innovation across Europe’s food and agriculture sector
\textbf{与你有关：} 命中用户订阅版面：world、business、ai；来源接近一手且权威度高；标题/正文包含高信号变化
{\footnotesize\color{ormuted} 来源：European Commission Digital Strategy News · \url{https://digital-strategy.ec.europa.eu/en/news/edic-agri-food-launched-support-digital-innovation-across-europes-food-and-agriculture-sector}}
\vspace{2mm}
\hrule height 0.2pt
\vspace{2mm}
{\large\bfseries Commission holds special meeting of Scientific panel on frontier AI safety and risks}\par
The Commission holds a special meeting of the Scientific Panel on artificial intelligence (AI). The panel has been investigating recent loss-of-control incidents, and together with the Commission’s AI Office, has worked on a set of questions for the companies that developed the models involved. Henna Virkkunen, Executive Vice-President for Tech Sovereignty, Security and Democracy, will attend the meeting. The Scienti
\textbf{与你有关：} 命中用户订阅版面：world、business、ai；来源接近一手且权威度高；标题/正文包含高信号变化
{\footnotesize\color{ormuted} 来源：European Commission Digital Strategy News · \url{https://digital-strategy.ec.europa.eu/en/news/commission-holds-special-meeting-scientific-panel-frontier-ai-safety-and-risks}}
\end{multicols}
\newpage
{\small\bfseries OR MORNING / SECTION 08 / 2026-10-13}
\vspace{2mm}
\hrule height 0.9pt
\vspace{1.5mm}
\hrule height 0.25pt
\vspace{4mm}
\section*{世界状态}
{\color{ormuted} 重大政策、国际、宏观与不应因为沉浸在 AI 中而错过的变化。}
\vspace{2mm}
{\LARGE\bfseries 行政机构\_中国人民大学高瓴人工智能学院}\par
\vspace{1mm}
{\color{ormuted} 本版筛出 5 个高信号事件。}
\vspace{3mm}
{\Large\bfseries 行政机构\_中国人民大学高瓴人工智能学院}\par
首页 学院概况 学院简介 数说学院 现任领导 学术委员会 行政机构 新闻公告 学院新闻 学院公告 师生风采 光影视界 师资队伍 专任教师 博士后队伍 人才招聘 人才培养 本科生 学术型硕士研究生 专业学位硕士研究生 学术型博士研究生 专业学位博士研究生 在职课程培训班 同等学力申请硕士学位 科学研究 研究机构 研究成果 合作交流 对外合作 国际交流 校友名录 光影记录 校友名录 高礼研究院 EN 教师入口 学院概况 学院简介 数说学院 现任领导 学术委员会 行政机构 行政机构 您当前所在的位置： 首页 > 学院概况 > 行政机构 教务工作 - 教务管理：负责排课、选课、成绩管理，以及毕业实践和毕业论文的组织与管理，包括出具成绩单和协助学历研究生招生工作。 - 教师教学工作管理：计算和考核教师的教学工作量，确保教学任务的合理分配和履行。 - 教学委员会支持：协助各教学工作委员会的秘书处工作，包括会议组织、文件管理等。 - 学科建设
\textbf{与你有关：} 这是当前用户可实际参与或申请的高可触达机会；参与范围：北京 / 人大；命中用户订阅版面：opportunity、research、ai；来源接近一手且权威度高；标题/正文包含高信号变化
\textbf{建议：} 检查报名资格、截止时间、参与形式和预期收益；合适则立即加入任务系统。
{\footnotesize\color{ormuted} 来源：中国人民大学高瓴人工智能学院公告 · \url{https://ai.ruc.edu.cn/xygk/xzjg/index.htm}}
\vspace{4mm}
\hrule height 0.25pt
\vspace{3mm}
\begin{multicols}{2}
{\large\bfseries 姚班培养-清华大学交叉信息研究院}\par
首页 热点新闻 讲座预告 院系概况 院系简介 院长寄语 现任领导 院系沿革 新闻动态 院系动态 媒体聚焦 人员队伍 全职教师 研究系列 荣誉客座 兼职教师 博士后 行政人员 实验员 姚班本科生 姚班培养 姚班招生 研究生 研究生培养 研究生招生 量子中心 中心简介 师资力量 科研概况 量子中心动态 科学研究 组织架构 课题组介绍 研究动态 科研简报 校园生活 学生活动 衣食住行 公开招聘 教师招聘 博士后招聘 职员招聘 校友频道 校友风采 校友俱乐部 校友捐赠 姚班本科生 姚班培养 姚班招生 清华学堂计算机科学实验班 姚班介绍 “清华学堂计算机科学实验班（姚班）”由世界著名计算机科学家姚期智院士于2005年创办，是清华大学培养国际顶尖计算机科学及交叉创新人才的本科生实验班，被誉为“最优秀的本科生和最优秀的本科教育”标杆。 姚班面向科技发展前沿与国家重大需求构建培养体系，以计算机科学、人工智能和量子信息三大前沿专业方向为核心，供学
\textbf{与你有关：} 这是当前用户可实际参与或申请的高可触达机会；参与范围：北京 / 清华；命中用户订阅版面：opportunity、research、ai；来源接近一手且权威度高；标题/正文包含高信号变化
\textbf{建议：} 检查报名资格、截止时间、参与形式和预期收益；合适则立即加入任务系统。
{\footnotesize\color{ormuted} 来源：清华交叉信息研究院 · \url{https://iiis.tsinghua.edu.cn/ybbks/ybpy.htm}}
\vspace{2mm}
\hrule height 0.2pt
\vspace{2mm}
{\large\bfseries 中心简介-清华大学交叉信息研究院}\par
首页 热点新闻 讲座预告 院系概况 院系简介 院长寄语 现任领导 院系沿革 新闻动态 院系动态 媒体聚焦 人员队伍 全职教师 研究系列 荣誉客座 兼职教师 博士后 行政人员 实验员 姚班本科生 姚班培养 姚班招生 研究生 研究生培养 研究生招生 量子中心 中心简介 师资力量 科研概况 量子中心动态 科学研究 组织架构 课题组介绍 研究动态 科研简报 校园生活 学生活动 衣食住行 公开招聘 教师招聘 博士后招聘 职员招聘 校友频道 校友风采 校友俱乐部 校友捐赠 量子中心 中心简介 师资力量 科研概况 量子中心动态 量子中心 中心介绍 量子信息是一门新兴的交叉学科领域，涉及计算机科学、信息科学和量子物理，为后摩尔时代提供了新的技术路径，探索了21世纪信息科学发展的新基础与新方法。量子信息技术已成为各国战略竞争的焦点之一。量子信息科学作为量子调控的核心，是信息网络与纳米科学等领域的交叉方向，代表了中国未来科技发展的重要战略方向之一。
\textbf{与你有关：} 这是当前用户可实际参与或申请的高可触达机会；参与范围：北京 / 清华；命中用户订阅版面：opportunity、research、ai；来源接近一手且权威度高；标题/正文包含高信号变化
\textbf{建议：} 检查报名资格、截止时间、参与形式和预期收益；合适则立即加入任务系统。
{\footnotesize\color{ormuted} 来源：清华交叉信息研究院 · \url{https://iiis.tsinghua.edu.cn/lzzx/zxjj.htm}}
\vspace{2mm}
\hrule height 0.2pt
\vspace{2mm}
{\large\bfseries 最新日程：清华大学第二十一届公共管理前沿论坛即将开幕}\par
最新日程：清华大学第二十一届公共管理前沿论坛即将开幕 IR AI SoftEng 清华大学人工智能国际治理研究院 2026-10-10 22:50 分享 以下文章来源于mp.weixin.qq.com 点击蓝字 关注我们 清华大学第二十一届公共管理前沿论坛将于2026年10月11日在清华大学正式举行。本届前沿论坛主题围绕“数智时代的国家治理现代化”，并就人工智能、自主知识体系、新兴技术发展、“十五五”规划、碳双控等议题设置六大平行论坛进行探讨和交流。本届论坛邀请了来自全国各地的学术界与实务界代表参会，欢迎广大院友、师生参会！ 活动议程 上午：主论坛 签到 时 间： 10月11日（周日） 8:00-8:30 地 点： 清华大学蒙民伟音乐厅 开幕式（8:30-9:00） 主持人： 谢 矜 清华大学公共管理学院党委书记 致辞嘉宾： 杨 斌 清华大学校务委员会副主任，清华大学经济管理学院教授 路易斯·瓦熙（Luis Vassy） 巴黎政
\textbf{与你有关：} 这是当前用户可实际参与或申请的高可触达机会；参与范围：北京 / 线上；命中用户订阅版面：opportunity、ai、research；来源接近一手且权威度高；标题/正文包含高信号变化
\textbf{建议：} 检查报名资格、截止时间、参与形式和预期收益；合适则立即加入任务系统。
{\footnotesize\color{ormuted} 来源：北京智源社区活动 · \url{https://hub.baai.ac.cn/view/58571}}
\vspace{2mm}
\hrule height 0.2pt
\vspace{2mm}
{\large\bfseries Back to Careers}\par
Back to Careers Join our team All Teams AI Research \& Engineering Applied AI Communications Compute Data Science \& Analytics Engineering \& Design - Product Finance Hardware Legal Marketing \& Brand People Product Management, Support, \& Operations Public Benefit Public Policy Safeguards (Trust \& Safety) Sales Security Software Engineering - Infrastructure Technical Education Technical Program Management All Offices Ban
\textbf{与你有关：} 参与范围：全球 / 多为海外全职；命中用户订阅版面：opportunity、ai、research；来源接近一手且权威度高；标题/正文包含高信号变化
\textbf{建议：} 检查报名资格、截止时间、参与形式和预期收益；合适则立即加入任务系统。
{\footnotesize\color{ormuted} 来源：Anthropic Careers · \url{https://www.anthropic.com/careers}}
\end{multicols}
\end{document}
